% Options for packages loaded elsewhere
\PassOptionsToPackage{unicode}{hyperref}
\PassOptionsToPackage{hyphens}{url}
\documentclass[
  a4paper,
  10pt]{article}
\usepackage{xcolor}
\usepackage[margin=22mm]{geometry}
\usepackage{amsmath,amssymb}
\setcounter{secnumdepth}{-\maxdimen} % remove section numbering
\usepackage{iftex}
\ifPDFTeX
  \usepackage[T1]{fontenc}
  \usepackage[utf8]{inputenc}
  \usepackage{textcomp} % provide euro and other symbols
\else % if luatex or xetex
  \usepackage{unicode-math} % this also loads fontspec
  \defaultfontfeatures{Scale=MatchLowercase}
  \defaultfontfeatures[\rmfamily]{Ligatures=TeX,Scale=1}
\fi
\usepackage{lmodern}
\ifPDFTeX\else
  % xetex/luatex font selection
\fi
% Use upquote if available, for straight quotes in verbatim environments
\IfFileExists{upquote.sty}{\usepackage{upquote}}{}
\IfFileExists{microtype.sty}{% use microtype if available
  \usepackage[]{microtype}
  \UseMicrotypeSet[protrusion]{basicmath} % disable protrusion for tt fonts
}{}
\makeatletter
\@ifundefined{KOMAClassName}{% if non-KOMA class
  \IfFileExists{parskip.sty}{%
    \usepackage{parskip}
  }{% else
    \setlength{\parindent}{0pt}
    \setlength{\parskip}{6pt plus 2pt minus 1pt}}
}{% if KOMA class
  \KOMAoptions{parskip=half}}
\makeatother
\usepackage{color}
\usepackage{fancyvrb}
\newcommand{\VerbBar}{|}
\newcommand{\VERB}{\Verb[commandchars=\\\{\}]}
\DefineVerbatimEnvironment{Highlighting}{Verbatim}{commandchars=\\\{\}}
% Add ',fontsize=\small' for more characters per line
\newenvironment{Shaded}{}{}
\newcommand{\AlertTok}[1]{\textcolor[rgb]{1.00,0.00,0.00}{\textbf{#1}}}
\newcommand{\AnnotationTok}[1]{\textcolor[rgb]{0.38,0.63,0.69}{\textbf{\textit{#1}}}}
\newcommand{\AttributeTok}[1]{\textcolor[rgb]{0.49,0.56,0.16}{#1}}
\newcommand{\BaseNTok}[1]{\textcolor[rgb]{0.25,0.63,0.44}{#1}}
\newcommand{\BuiltInTok}[1]{\textcolor[rgb]{0.00,0.50,0.00}{#1}}
\newcommand{\CharTok}[1]{\textcolor[rgb]{0.25,0.44,0.63}{#1}}
\newcommand{\CommentTok}[1]{\textcolor[rgb]{0.38,0.63,0.69}{\textit{#1}}}
\newcommand{\CommentVarTok}[1]{\textcolor[rgb]{0.38,0.63,0.69}{\textbf{\textit{#1}}}}
\newcommand{\ConstantTok}[1]{\textcolor[rgb]{0.53,0.00,0.00}{#1}}
\newcommand{\ControlFlowTok}[1]{\textcolor[rgb]{0.00,0.44,0.13}{\textbf{#1}}}
\newcommand{\DataTypeTok}[1]{\textcolor[rgb]{0.56,0.13,0.00}{#1}}
\newcommand{\DecValTok}[1]{\textcolor[rgb]{0.25,0.63,0.44}{#1}}
\newcommand{\DocumentationTok}[1]{\textcolor[rgb]{0.73,0.13,0.13}{\textit{#1}}}
\newcommand{\ErrorTok}[1]{\textcolor[rgb]{1.00,0.00,0.00}{\textbf{#1}}}
\newcommand{\ExtensionTok}[1]{#1}
\newcommand{\FloatTok}[1]{\textcolor[rgb]{0.25,0.63,0.44}{#1}}
\newcommand{\FunctionTok}[1]{\textcolor[rgb]{0.02,0.16,0.49}{#1}}
\newcommand{\ImportTok}[1]{\textcolor[rgb]{0.00,0.50,0.00}{\textbf{#1}}}
\newcommand{\InformationTok}[1]{\textcolor[rgb]{0.38,0.63,0.69}{\textbf{\textit{#1}}}}
\newcommand{\KeywordTok}[1]{\textcolor[rgb]{0.00,0.44,0.13}{\textbf{#1}}}
\newcommand{\NormalTok}[1]{#1}
\newcommand{\OperatorTok}[1]{\textcolor[rgb]{0.40,0.40,0.40}{#1}}
\newcommand{\OtherTok}[1]{\textcolor[rgb]{0.00,0.44,0.13}{#1}}
\newcommand{\PreprocessorTok}[1]{\textcolor[rgb]{0.74,0.48,0.00}{#1}}
\newcommand{\RegionMarkerTok}[1]{#1}
\newcommand{\SpecialCharTok}[1]{\textcolor[rgb]{0.25,0.44,0.63}{#1}}
\newcommand{\SpecialStringTok}[1]{\textcolor[rgb]{0.73,0.40,0.53}{#1}}
\newcommand{\StringTok}[1]{\textcolor[rgb]{0.25,0.44,0.63}{#1}}
\newcommand{\VariableTok}[1]{\textcolor[rgb]{0.10,0.09,0.49}{#1}}
\newcommand{\VerbatimStringTok}[1]{\textcolor[rgb]{0.25,0.44,0.63}{#1}}
\newcommand{\WarningTok}[1]{\textcolor[rgb]{0.38,0.63,0.69}{\textbf{\textit{#1}}}}
\usepackage{longtable,booktabs,array}
\newcounter{none} % for unnumbered tables
\usepackage{calc} % for calculating minipage widths
% Correct order of tables after \paragraph or \subparagraph
\usepackage{etoolbox}
\makeatletter
\patchcmd\longtable{\par}{\if@noskipsec\mbox{}\fi\par}{}{}
\makeatother
% Allow footnotes in longtable head/foot
\IfFileExists{footnotehyper.sty}{\usepackage{footnotehyper}}{\usepackage{footnote}}
\makesavenoteenv{longtable}
\setlength{\emergencystretch}{3em} % prevent overfull lines
\providecommand{\tightlist}{%
  \setlength{\itemsep}{0pt}\setlength{\parskip}{0pt}}
\usepackage{bookmark}
\IfFileExists{xurl.sty}{\usepackage{xurl}}{} % add URL line breaks if available
\urlstyle{same}
\hypersetup{
  hidelinks,
  pdfcreator={LaTeX via pandoc}}

\author{}
\date{}

\begin{document}

\section{Emergence of Discrete Born-Type Weights in Iterated Two-Channel
Exchange
Systems}\label{emergence-of-discrete-born-type-weights-in-iterated-two-channel-exchange-systems}

\subsection{A Finite-Order Recurrence Law from Wave-Packet Localization
Transfer, Metastable Two-State Dynamics, and Observation
Selection}\label{a-finite-order-recurrence-law-from-wave-packet-localization-transfer-metastable-two-state-dynamics-and-observation-selection}

\textbf{Version:} English complete manuscript v1 \textbf{Date:} July 18,
2026 \textbf{Author:} Noriaki Kihara \textbf{ORCID:} 0009-0004-6753-4020
\textbf{Version DOI:} 10.5281/zenodo.21422471 \textbf{Concept DOI:}
10.5281/zenodo.21422470 \textbf{Position in the series:} Independent
report on finite-order recurrence and Born-type weights in the Wave
Information Readout series

\begin{center}\rule{0.5\linewidth}{0.5pt}\end{center}

\subsection{Abstract}\label{abstract}

This study was not initiated to reproduce the Born rule or its two-state
representation, the \(\cos^2\) law. It began with two different
numerical experiments. System A examined whether iterative exchange
scattering between a weakly localized wave and a localized wave
containing higher harmonics could transfer localization and effective
harmonic structure between two channels, producing behavior resembling
wave-packet contraction or localization. System B examined two A/B
states, an intermediate gray metastable state, its retention under weak
readout, and its selection into A or B under strong observation.

Although the two models were built to address different physical
questions and used different observation functions, later analysis
showed that their linear mixing parts, before normalization and
observation operations, share the same two-channel exchange-scattering
kernel

\[
U_R=
\begin{pmatrix}
r&t\\
t&r
\end{pmatrix}.
\]

System A applies channel-wise normalization after this linear mixing,
while the strong D observation in System B applies a separate
back-action map; the complete update maps of the two systems are
therefore not assumed to be identical. A full-range sweep of the
exchange coefficient \(R\) in System B produced extremely sharp peaks at
particular values. Eigenvalue analysis showed that the peaks arise from
exact recurrence in which the antisymmetric eigenvalue returns to unity
after finitely many iterations:

\[
U_R^n=I.
\]

We derive this result without depending on the trigonometric
parameterization of the scattering amplitudes. Exchange symmetry
determines the symmetric and antisymmetric projectors

\[
P_s=\frac12
\begin{pmatrix}
1&1\\
1&1
\end{pmatrix},
\qquad
P_a=\frac12
\begin{pmatrix}
1&-1\\
-1&1
\end{pmatrix}.
\]

After fixing the global phase so that the symmetric eigenvalue is one,
an exchange-symmetric unitary operator can be written as

\[
U=P_s+\zeta P_a,
\qquad |\zeta|=1.
\]

When the finite-order condition selects \(\zeta=e^{-2\pi im/n}\), the
diagonal and off-diagonal amplitudes in the A/B basis are

\[
r=\frac{1+\zeta}{2},
\qquad
t=\frac{1-\zeta}{2}.
\]

Consequently, for coprime integers \(m,n\), the exchange weight
satisfying the finite-order condition is

\[
\boxed{
R_{n,m}=\cos^2\left(\frac{\pi m}{n}\right)
}.
\]

This expression was not obtained by entering a Born-type squared law as
a search condition or evaluation function. The projector representation
above also shows that the trigonometric form need not be imposed as an
independent assumption. The iterative exchange systems modeling
wave-packet localization transfer, metastable two-state dynamics, weak
readout, and strong observation selection were constructed first;
numerical peaks were then found; only afterward was their closed-orbit
condition analyzed. The result is therefore not a complete derivation of
the Born rule. It shows that, in a closed two-channel wave system
exhibiting behavior analogous to phenomena associated with quantum
measurement, finitely recurrent exchange weights arise endogenously as a
discrete Born-type \(\cos^2\) series.

The current code represents the unitary scattering amplitudes using
\(\sin\theta\) and \(\cos\theta\). This is a coordinate representation
recoverable from the exchange-symmetric projectors \(P_s,P_a\) and the
eigenvalues \(1,\zeta\), rather than an independent assumption required
by the central theorem. No angle was prescribed in advance as the one
that must recur exactly. Instead, the finite-order condition selects the
discrete Born-angle series

\[
\phi_{n,m}=\frac{\pi m}{n}.
\]

The scattering angle used in the implementation is
\(\theta=\pi/2-\phi_{n,m}\). This distinguishes the exchange-symmetric
two-channel projection structure from the phase series selected by
closure: the latter discretizes the admissible weights of the former.

This paper does not focus on numerical correspondences with the
fine-structure constant, \(E_8\), or any specific physical constant. Its
central purpose is to report clearly the mathematical structure
connecting contraction-like localization redistribution, metastable
superposition, observation selection, finite-order recurrence, and
Born-type squared weights through a common unperturbed exchange kernel
\(U\) followed by model-specific normalization, readout, and selection
maps.

\textbf{Keywords:} Born rule, two-channel exchange, finite-order
recurrence, wave-packet localization, metastable state, observation
selection, quantum measurement, unitary operator, discrete phase

\begin{center}\rule{0.5\linewidth}{0.5pt}\end{center}

\section{1. Background and Objective}\label{1-background-and-objective}

\subsection{1.1 The Born Rule and Its Two-State
Representation}\label{11-the-born-rule-and-its-two-state-representation}

In standard quantum mechanics, the probability of projecting a state
\(|\psi\rangle\) onto an observation state \(|A\rangle\) is given by the
Born rule

\[
P(A)=|\langle A|\psi\rangle|^2
\]

{[}5,6{]}.

For a two-state system written as

\section{\$\$
\textbar\textbackslash psi\textbackslash rangle}\label{-psirangle}

\textbackslash cos\textbackslash phi,\textbar A\textbackslash rangle +
e\^{}\{i\textbackslash chi\}\textbackslash sin\textbackslash phi,\textbar B\textbackslash rangle,
\$\$

the observation probabilities of the two orthogonal states are

\[
P(A)=\cos^2\phi,
\qquad
P(B)=\sin^2\phi.
\]

The Born rule is the fundamental rule connecting quantum theory to
observation results. The present work, however, did not begin as a
derivation of that rule. The numerical systems were first constructed to
study wave-packet localization transfer and observation selection in a
metastable two-state system; they did not contain the \(\cos^2\) law as
a target function.

\subsection{1.2 Research Question}\label{12-research-question}

The question addressed here is:

\begin{quote}
In an iterated exchange system that models contraction-like localization
redistribution, metastable mixing of two A/B states, retention under
weak readout, and selection under strong observation, can Born-type
squared weights arise from exchange symmetry, unitarity, and
finite-order closure without being externally entered as a probability
axiom?
\end{quote}

This paper does not seek a uniqueness proof of the general Born rule or
a complete dynamical derivation of single-trial probabilities. Its
specific objective is to decompose an already constructed
exchange-symmetric two-channel unitary operator into symmetric and
antisymmetric subspaces, derive the endogenous discrete Born-type
weights

\[
R_{n,m}=\cos^2\left(\frac{\pi m}{n}\right)
\]

from its finite-order condition, and state precisely the meaning and
limits of this result.

\subsection{1.3 The Non-Directed Order of
Discovery}\label{13-the-non-directed-order-of-discovery}

The chronological order of the discovery is important to its
interpretation.

\begin{enumerate}
\def\labelenumi{\arabic{enumi}.}
\tightlist
\item
  System A was constructed to transfer wave-packet localization and
  harmonic structure between two channels {[}1,2{]}.
\item
  System B was constructed to study two A/B states, a gray metastable
  state, weak readout, and strong observation selection {[}3{]}.
\item
  A local sweep was performed for localization transfer in System A, and
  a full-range sweep of their common exchange-scattering coefficient
  \(R\) was performed in System B {[}4{]}.
\item
  Extremely sharp peaks were observed in the full-range System B sweep.
\item
  Local high-resolution sweeps and multiprecision calculations were
  performed.
\item
  Operator eigenvalues were analyzed, identifying the peaks as
  finite-order recurrence {[}4{]}.
\item
  The finite-order condition was recognized to give the exchange weights
  \(\cos^2(\pi m/n)\).
\item
  Only afterward was this expression reinterpreted physically as
  isomorphic to the two-state representation of the Born rule.
\end{enumerate}

Thus, this study did not embed a Born-type \(\cos^2\) law in the code
and rediscover it in the output. The amplitude structure of unitary
two-channel scattering was present in the original models, but neither
the phase angle nor the exchange weight selected for exact recurrence
was entered in advance. Moreover, this paper removes the trigonometric
representation itself from the central assumptions and derives

\[
U=P_s+\zeta P_a,
\qquad |\zeta|=1
\]

from exchange symmetry and unitarity. Imposing the finite-order
condition on this representation makes the \(\cos^2/\sin^2\) weights a
necessary consequence of interference between the symmetric and
antisymmetric eigenphases, rather than of a chosen coordinate
parameterization.

\begin{center}\rule{0.5\linewidth}{0.5pt}\end{center}

\section{2. The Two Starting Models}\label{2-the-two-starting-models}

\subsection{2.1 System A: Wave-Packet Localization
Transfer}\label{21-system-a-wave-packet-localization-transfer}

System A iterates the same exchange scattering between a weakly
localized wave on side A and a localized wave containing higher
harmonics on side B {[}1,2{]}.

Let the wave-packet vectors at step \(j\) be \(A_j,B_j\). The linear
exchange part is

\section{\$\$ \textbackslash begin\{pmatrix\} \textbackslash widetilde
A\_\{j+1\}\textbackslash{} \textbackslash widetilde B\_\{j+1\}
\textbackslash end\{pmatrix\}}\label{-beginpmatrix-widetilde-a_j1-widetilde-b_j1-endpmatrix}

U\_R \textbackslash begin\{pmatrix\} A\_j\textbackslash{} B\_j
\textbackslash end\{pmatrix\}. \$\$

The implemented one-step update in System A is not only this linear
mixing. If \(\mathcal N(x):=x/\|x\|\) denotes normalization of each
wave-packet channel, the actual update is

\section{\$\$ \textbackslash boxed\{ F\_R(A,B)}\label{-boxed-f_rab}

\textbackslash left( \textbackslash mathcal N(rA+tB),
\textbackslash mathcal N(tA+rB) \textbackslash right) \}. \$\$

Thus, \(F_R\) is generally nonlinear. The finite-order theorem in this
paper applies directly to the pre-normalization linear numerator

\[
(A,B)\longmapsto(rA+tB,\ tA+rB).
\]

The localization transfer observed in System A is treated as a property
of the complete map combining this common linear kernel with
channel-wise normalization.

Representative observables are

\[
L,
\qquad
N_{\mathrm{eff}},
\qquad
B_{\mathrm{to}A},
\]

where

\begin{itemize}
\tightlist
\item
  \(L\) measures wave-packet localization,
\item
  \(N_{\mathrm{eff}}\) is the effective harmonic order, and
\item
  \(B_{\mathrm{to}A}\) measures how much of the initial harmonic
  structure on side B has moved to side A.
\end{itemize}

At intermediate exchange coefficients, localization and harmonic
structure transfer periodically between the two channels. If only one
channel is observed at a selected time, a broad wave packet can appear
to have changed into a localized wave. This behavior was constructed as
redistribution of localization by exchange interference, not as
nonunitary contraction accompanied by disappearance of the norm of the
total system.

\subsection{2.2 System B: White-Cat, Black-Cat, and Gray-Cat Metastable
Interface}\label{22-system-b-white-cat-black-cat-and-gray-cat-metastable-interface}

System B defines, from two complex amplitudes \(a,b\),

\[
p_A=|a|^2,
\qquad
p_B=|b|^2,
\qquad
S=p_A-p_B
\]

{[}3{]}.

The states are read as follows.

\begin{itemize}
\tightlist
\item
  \(S\approx+1\): A-dominant, or white-cat, state.
\item
  \(S\approx-1\): B-dominant, or black-cat, state.
\item
  \(S\approx0\): balanced A/B, or gray, state.
\end{itemize}

Among gray states, a state that fluctuates periodically with small
amplitude, remains under weak C readout, and is selected into A or B by
strong D observation is called a gray metastable phase.

The unperturbed A/B exchange in System B follows \(U_R\). The
implemented pair normalization is

\section{\$\$ \textbackslash mathcal
N\_\{AB\}(a,b)}\label{-mathcal-n_abab}

\textbackslash frac\{(a,b)\}\{\textbackslash sqrt\{\textbar a\textbar\^{}2+\textbar b\textbar\^{}2\}\},
\$\$

which is the identity under exact unitary evolution. The weak C readout
and strong D observation are separate readout and back-action maps
applied after the unperturbed exchange. In particular, the D
implementation explicitly contains nonlinear back-action that updates
the allocation difference \(S\) as

\section{\$\$ S\_\{j+1\}}\label{-s_j1}

S\_j+g\_D S\_\{D,j\}(1-S\_j\^{}2), \$\$

pushing the state toward either A or B. Strong observation selection
therefore does not arise from attraction by \(U_R\) alone, but from the
composition of \(U_R\) with D back-action.

This model was not constructed as a direct solution to the
quantum-measurement problem. Its purpose was to test whether one
iterative exchange system could distinguish

\begin{itemize}
\tightlist
\item
  a metastable state intermediate between two states,
\item
  a weak readout that does not substantially destroy the state, and
\item
  a strong observation that changes the state toward one side.
\end{itemize}

\subsection{2.3 The Scattering Kernel Shared by the Two
Models}\label{23-the-scattering-kernel-shared-by-the-two-models}

Later analysis showed that System A and System B, despite their
different state spaces, normalizations, and observation functions, use
the same exchange-scattering kernel in the linear mixing stage before
normalization, readout, and back-action:

\[
\boxed{
U_R=
\begin{pmatrix}
r&t\\
t&r
\end{pmatrix}
}.
\]

The complete update structures can be decomposed conceptually as

\[
\text{System A}:\quad F_R=\mathcal N_A\circ U_R,
\]

\[
\text{System B, unperturbed}:\quad U_R,
\]

\[
\text{System B, C/D}:\quad
F_C=\mathcal C\circ U_R,
\qquad
F_D=\mathcal D\circ U_R.
\]

Here \(\mathcal N_A\) is channel-wise normalization, \(\mathcal C\) is
weak readout and weak back-action, and \(\mathcal D\) is strong
selection back-action. System A reads localization and internal harmonic
structure of wave-packet vectors, whereas System B reads the A/B
allocation difference and metastability of two complex amplitudes.

Therefore, candidate positions of particular \(R\) values common to both
models originate in the intrinsic structure of the exchange operator,
while peak visibility, localization transfer, and A/B selection are
determined by the downstream maps and observation functions. This
separation between the common kernel that fixes root positions and the
model-specific maps that visualize or select phenomena is the central
implementation structure of this paper.

\begin{center}\rule{0.5\linewidth}{0.5pt}\end{center}

\section{3. Exchange-Scattering
Operator}\label{3-exchange-scattering-operator}

\subsection{3.1 Scattering-Coefficient Coordinates Used in the Numerical
Implementation}\label{31-scattering-coefficient-coordinates-used-in-the-numerical-implementation}

For a reflection weight \(R\in[0,1]\), the implementation defines

\[
\theta(R):=\arcsin\sqrt R.
\]

The transmission amplitude \(t\) and reflection amplitude \(r\) are

\[
t=e^{i\theta}\cos\theta,
\qquad
r=-ie^{i\theta}\sin\theta.
\]

Then

\[
|t|^2=\cos^2\theta=1-R,
\]

\[
|r|^2=\sin^2\theta=R.
\]

Moreover,

\[
|r|^2+|t|^2=1,
\qquad
r^*t+t^*r=0,
\]

so that

\[
U_R^\dagger U_R=I.
\]

Thus, the linear exchange operator at fixed \(R\) is unitary. This
trigonometric representation is a convenient coordinate system used in
the implementation, but it is not an independent axiom needed to obtain
the \(\cos^2/\sin^2\) series. Sections 4 and 5 reconstruct the same
coefficient representation using only the projector decomposition of an
exchange-symmetric unitary operator.

\subsection{3.2 State Norm and
Localization}\label{32-state-norm-and-localization}

The unitarity of the linear kernel \(U_R\) preserves the norm of the
complete two-channel state:

\section{\$\$ \textbackslash left\textbar{} U\_R
\textbackslash begin\{pmatrix\} A\textbackslash B
\textbackslash end\{pmatrix\}
\textbackslash right\textbar\^{}2}\label{-left-u_r-beginpmatrix-ab-endpmatrix-right2}

\textbackslash left\textbar{} \textbackslash begin\{pmatrix\}
A\textbackslash B \textbackslash end\{pmatrix\}
\textbackslash right\textbar\^{}2. \$\$

Therefore, even if one wave packet appears to localize under the
pre-normalization linear kernel, this does not mean that the total norm
disappears or that the state is dissipatively attracted to a point.
Localization, harmonic structure, or channel allocation is redistributed
within the conserved total state. Since the complete System A map
includes channel-wise normalization, this total-norm equation cannot be
transferred unchanged to the full update; localization transfer must be
evaluated for \(F_R=\mathcal N_A\circ U_R\).

This distinction is essential when discussing the analogy with
wave-packet contraction. ``Contraction-like behavior'' in this paper is
not claimed to be identical to post-measurement state update in standard
quantum mechanics. It denotes a mathematical analogy in which
concentration of localization emerges from the composition of a unitary
exchange kernel, channel-wise normalization, and partial observation.

\begin{center}\rule{0.5\linewidth}{0.5pt}\end{center}

\section{4. Symmetric and Antisymmetric
Eigenmodes}\label{4-symmetric-and-antisymmetric-eigenmodes}

\subsection{4.1 Decomposition into
Eigenmodes}\label{41-decomposition-into-eigenmodes}

Decompose the two-channel state as

\[
X_j:=\frac{A_j+B_j}{\sqrt2},
\qquad
Y_j:=\frac{A_j-B_j}{\sqrt2}.
\]

\(X_j\) is symmetric under channel exchange, and \(Y_j\) is
antisymmetric.

The eigenvalues of the exchange operator are

\[
\lambda_s=r+t,
\qquad
\lambda_a=r-t.
\]

Substituting the definitions of \(r,t\) gives

\[
\lambda_s=1,
\]

\[
\lambda_a=-e^{2i\theta}.
\]

Therefore,

\[
X_j=X_0,
\qquad
Y_j=\lambda_a^jY_0.
\]

\subsection{4.2 Projector
Representation}\label{42-projector-representation}

Let the channel-exchange operator be

\[
X=
\begin{pmatrix}
0&1\\
1&0
\end{pmatrix}
=P_s-P_a.
\]

Exchange symmetry is

\[
[U,X]=0.
\]

Because the eigenvalues \(+1,-1\) of \(X\) are distinct, any operator
\(U\) satisfying this commutation relation does not mix the symmetric
and antisymmetric subspaces.

Define the symmetric and antisymmetric projectors by

\[
P_s=
\frac12
\begin{pmatrix}
1&1\\
1&1
\end{pmatrix},
\qquad
P_a=
\frac12
\begin{pmatrix}
1&-1\\
-1&1
\end{pmatrix}.
\]

Then

\[
P_s+P_a=I,
\qquad
P_sP_a=0,
\]

and

\[
U=\lambda_sP_s+\lambda_aP_a.
\]

If \(U\) is unitary,

\[
|\lambda_s|=|\lambda_a|=1.
\]

Removing the global phase \(\lambda_s\) by defining
\(\widetilde U:=\lambda_s^{-1}U\) gives

\[
\boxed{
\widetilde U=P_s+\zeta P_a,
\qquad
\zeta:=\frac{\lambda_a}{\lambda_s},
\qquad
|\zeta|=1
}.
\]

In the numerical implementation \(\lambda_s=1\) from the outset, so
below we write \(\widetilde U\) as \(U_R\). Then

\section{\$\$ U\_R}\label{-u_r}

\section{\textbackslash frac12 \textbackslash begin\{pmatrix\}
1+\textbackslash zeta\&1-\textbackslash zeta\textbackslash{}
1-\textbackslash zeta\&1+\textbackslash zeta
\textbackslash end\{pmatrix\}}\label{frac12-beginpmatrix-1zeta1-zeta-1-zeta1zeta-endpmatrix}

\textbackslash begin\{pmatrix\} r\&t\textbackslash{} t\&r
\textbackslash end\{pmatrix\}, \$\$

and hence

\[
\boxed{
r=\frac{1+\zeta}{2},
\qquad
t=\frac{1-\zeta}{2}
}.
\]

Thus, without separately assuming a trigonometric form, the exchange
amplitudes are fixed as the sum and difference of the symmetric
eigenphase \(1\) and the antisymmetric eigenphase \(\zeta\).

Furthermore,

\[
\boxed{
U_R^j=P_s+\zeta^jP_a
}.
\]

This expression shows that the apparently complicated time evolution of
the model consists of only

\begin{itemize}
\tightlist
\item
  a fixed symmetric component, and
\item
  an antisymmetric component rotating on the unit circle.
\end{itemize}

\subsection{4.3 Periodicity of
Observables}\label{43-periodicity-of-observables}

The allocation difference in System B can generally be written as

\[
S_j=C\cos(j\omega+\varphi_0),
\]

where

\[
\omega=\arg\lambda_a
\]

is determined by the operator, while \(C\) and \(\varphi_0\) are
determined by the initial state and observation function.

Changing the initial amplitude, initial phase, internal harmonic
distribution, or wave-packet shape therefore does not change the
operator-intrinsic phase-rotation number \(\omega(R)\). What changes is
how visibly that eigenphase appears in the chosen observable.

\begin{center}\rule{0.5\linewidth}{0.5pt}\end{center}

\section{5. Finite-Order Recurrence}\label{5-finite-order-recurrence}

\subsection{5.1 Exact Recurrence
Condition}\label{51-exact-recurrence-condition}

In the numerical implementation, the symmetric eigenvalue is exactly
one. The condition under which this phase-fixed iterated exchange system
restores every initial state after \(n\) iterations is

\[
U_R^n=I.
\]

Since the symmetric eigenvalue satisfies

\[
\lambda_s^n=1,
\]

the necessary and sufficient condition is

\[
\lambda_a^n=1
\]

for the antisymmetric eigenvalue. For a general exchange-symmetric
\(U\in U(2)\), \(U^n=e^{i\gamma}I\) is projective recurrence up to a
global phase of the physical state. The equation \(\widetilde U^n=I\)
for \(\widetilde U=\lambda_s^{-1}U\) is its phase-fixed representation
used below.

\subsection{5.2 Exchange-Symmetric Finite-Order
Theorem}\label{52-exchange-symmetric-finite-order-theorem}

\textbf{Theorem 5.1 (Discrete squared weights of an exchange-symmetric
finite-order map).}\\
Let a two-channel operator \(U\in U(2)\) satisfy the following
conditions.

\begin{enumerate}
\def\labelenumi{\arabic{enumi}.}
\tightlist
\item
  It commutes with the channel-exchange operator \(X\): \([U,X]=0\).
\item
  The global phase is fixed by using \(\widetilde U:=\lambda_s^{-1}U\).
\item
  \(\widetilde U\) has a nontrivial fundamental order \(n\ge3\):
  \(\widetilde U^n=I\).
\item
  The antisymmetric eigenphase is \(\zeta=e^{-2\pi i m/n}\) with
  \(\gcd(m,n)=1\). Since the conjugate pair \(m\) and \(n-m\) gives the
  same weights, choose \(1\le m<n/2\) as a representative.
\end{enumerate}

Then

\section{\$\$ \textbackslash widetilde U}\label{-widetilde-u}

\section{P\_s+e\^{}\{-2\textbackslash pi i
m/n\}P\_a}\label{p_se-2pi-i-mnp_a}

\textbackslash begin\{pmatrix\} r\_\{n,m\}\&t\_\{n,m\}\textbackslash{}
t\_\{n,m\}\&r\_\{n,m\} \textbackslash end\{pmatrix\}, \$\$

\section{\$\$ r\_\{n,m\}}\label{-r_nm}

\section{\textbackslash frac\{1+e\^{}\{-2\textbackslash pi i
m/n\}\}\{2\}, \textbackslash qquad
t\_\{n,m\}}\label{frac1e-2pi-i-mn2-qquad-t_nm}

\textbackslash frac\{1-e\^{}\{-2\textbackslash pi i m/n\}\}\{2\}, \$\$

and the channel weights are

\section{\$\$ \textbackslash boxed\{
\textbar r\_\{n,m\}\textbar\^{}2}\label{-boxed-r_nm2}

\section{\textbackslash cos\^{}2\textbackslash left(\textbackslash frac\{\textbackslash pi
m\}\{n\}\textbackslash right), \textbackslash qquad
\textbar t\_\{n,m\}\textbar\^{}2}\label{cos2leftfracpi-mnright-qquad-t_nm2}

\textbackslash sin\^{}2\textbackslash left(\textbackslash frac\{\textbackslash pi
m\}\{n\}\textbackslash right) \}. \$\$

\emph{Proof.} The relation \([U,X]=0\) diagonalizes \(U\) on the
symmetric and antisymmetric eigenspaces of \(X\). After fixing the
global phase,

\[
\widetilde U=P_s+\zeta P_a.
\]

The condition \(\widetilde U^n=I\) gives \(\zeta^n=1\). Since the
fundamental order is \(n\), write \(\zeta=e^{-2\pi im/n}\) with
\(\gcd(m,n)=1\). Expanding the projectors in the channel basis gives

\[
r_{n,m}=\frac{1+\zeta}{2},
\qquad
t_{n,m}=\frac{1-\zeta}{2}.
\]

Therefore,

\section{\$\$ \textbar r\_\{n,m\}\textbar\^{}2}\label{-r_nm2}

\section{\textbackslash frac\{(1+\textbackslash zeta)(1+\textbackslash zeta\^{}*)\}\{4\}}\label{frac1zeta1zeta4}

\section{\textbackslash frac\{1+\textbackslash cos(2\textbackslash pi
m/n)\}\{2\}}\label{frac1cos2pi-mn2}

\textbackslash cos\^{}2\textbackslash left(\textbackslash frac\{\textbackslash pi
m\}\{n\}\textbackslash right), \$\$

\section{\$\$ \textbar t\_\{n,m\}\textbar\^{}2}\label{-t_nm2}

\section{\textbackslash frac\{1-\textbackslash cos(2\textbackslash pi
m/n)\}\{2\}}\label{frac1-cos2pi-mn2}

\textbackslash sin\^{}2\textbackslash left(\textbackslash frac\{\textbackslash pi
m\}\{n\}\textbackslash right). \textbackslash qquad\textbackslash square
\$\$

\subsection{5.3 Agreement with the Implementation Coordinates and
Boundary
Roots}\label{53-agreement-with-the-implementation-coordinates-and-boundary-roots}

In the implementation coordinates of Section 3,

\[
\zeta=\lambda_a=-e^{2i\theta}.
\]

Identifying this with \(e^{-2\pi im/n}\) gives

\section{\$\$ \textbackslash theta}\label{-theta}

\textbackslash frac\{\textbackslash pi\}\{2\}-\textbackslash frac\{\textbackslash pi
m\}\{n\} \textbackslash pmod\{\textbackslash pi\}, \$\$

and once again

\section{\$\$
R=\textbar r\textbar\^{}2=\textbackslash sin\^{}2\textbackslash theta}\label{-rr2sin2theta}

\textbackslash cos\^{}2\textbackslash left(\textbackslash frac\{\textbackslash pi
m\}\{n\}\textbackslash right). \$\$

The trigonometric form is therefore not an input to the theorem. It is
the implementation coordinate expressing the coefficients obtained from
the projector decomposition in terms of \(R\).

Theorem 5.1 concerns nontrivial interior roots \(0<R<1\). At the
boundaries,

\[
\zeta=1:\quad U=I,\quad R=1,
\]

\[
\zeta=-1:\quad U=X,\quad R=0,
\]

which are the trivial closures of order one and two, respectively.

\subsection{5.4 A Discrete Series inside a Continuous
Parameter}\label{54-a-discrete-series-inside-a-continuous-parameter}

Formally, \(R\) ranges over the continuous interval

\[
R\in[0,1].
\]

The points of exact recurrence, however, are classified by integer pairs
\((n,m)\):

\section{\$\$ (n,m) \textbackslash longmapsto
R\_\{n,m\}}\label{-nm-longmapsto-r_nm}

\textbackslash cos\^{}2\textbackslash left(\textbackslash frac\{\textbackslash pi
m\}\{n\}\textbackslash right). \$\$

Thus, within an apparently continuous exchange system there is a
countable family of closed orbits labeled by rational phase ratios.

Because rational numbers \(m/n\) are dense in the interval, finite-order
roots also become dense when arbitrarily high orders are allowed. At
each root, however,

\[
U_R^n=I
\]

holds exactly: the orbit is an exact discrete closed orbit, not an
approximately periodic one.

\begin{center}\rule{0.5\linewidth}{0.5pt}\end{center}

\section{6. Discrete Born-Type
Weights}\label{6-discrete-born-type-weights}

\subsection{6.1 Formal Isomorphism with Standard Two-State
Projection}\label{61-formal-isomorphism-with-standard-two-state-projection}

In a standard two-state system,

\section{\$\$
\textbar\textbackslash psi\textbackslash rangle}\label{-psirangle-1}

\textbackslash cos\textbackslash phi,\textbar A\textbackslash rangle +
e\^{}\{i\textbackslash chi\}\textbackslash sin\textbackslash phi,\textbar B\textbackslash rangle,
\$\$

the Born rule gives

\[
P(A)=\cos^2\phi,
\qquad
P(B)=\sin^2\phi.
\]

Applying the operator studied here to the channel-basis state
\(|A\rangle=(1,0)^T\) gives

\section{\$\$
U\_\{n,m\}\textbar A\textbackslash rangle}\label{-u_nmarangle}

r\_\{n,m\}\textbar A\textbackslash rangle+t\_\{n,m\}\textbar B\textbackslash rangle.
\$\$

Consequently, the two transition weights obtained by projecting the
output onto the same A/B basis are exactly

\section{\$\$ W\_\{A\textbackslash to A\}}\label{-w_ato-a}

\section{\textbar\textbackslash langle
A\textbar U\_\{n,m\}\textbar A\textbackslash rangle\textbar\^{}2}\label{langle-au_nmarangle2}

\textbar r\_\{n,m\}\textbar\^{}2, \$\$

\section{\$\$ W\_\{A\textbackslash to B\}}\label{-w_ato-b}

\section{\textbar\textbackslash langle
B\textbar U\_\{n,m\}\textbar A\textbackslash rangle\textbar\^{}2}\label{langle-bu_nmarangle2}

\textbar t\_\{n,m\}\textbar\^{}2. \$\$

By Theorem 5.1,

\section{\$\$ W\_\{A\textbackslash to A\}(n,m)}\label{-w_ato-anm}

\textbackslash cos\^{}2\textbackslash phi\_\{n,m\}, \$\$

\section{\$\$ W\_\{A\textbackslash to B\}(n,m)}\label{-w_ato-bnm}

\textbackslash sin\^{}2\textbackslash phi\_\{n,m\}, \$\$

where

\[
\phi_{n,m}:=\frac{\pi m}{n}.
\]

Thus, the channel-projection weights of a single exchange acting on a
basis input are

\[
\boxed{
W_{A\to A}(n,m)=\cos^2\phi_{n,m},
\qquad
W_{A\to B}(n,m)=\sin^2\phi_{n,m}
}
\]

and are isomorphic to the standard two-state Born rule. This is not
merely a visual resemblance between formulas; it is the same operation
of projecting a complex amplitude onto an observation basis and squaring
its absolute value.

For a general superposed input \(\alpha|A\rangle+\beta|B\rangle\),
however, the A-channel output is \(r\alpha+t\beta\), whose weight
includes an interference term. Therefore, \(R=|r|^2\) is not itself the
probability of observing A for an arbitrary input state; it is the
transition weight for remaining in the A channel when the input is the A
basis state. A separate frequency law is also required to interpret this
projection weight as an experimental frequency over repeated trials.

\subsection{6.2 What Was Assumed and What
Emerged}\label{62-what-was-assumed-and-what-emerged}

To evaluate the result precisely, we separate the structures already
present in the initial definitions from those that emerged only after
analysis.

\subsubsection{Structures present in the initial
definitions}\label{structures-present-in-the-initial-definitions}

\begin{enumerate}
\def\labelenumi{\arabic{enumi}.}
\tightlist
\item
  Two-channel complex amplitudes.
\item
  A unitary scattering condition.
\item
  Exchange symmetry between the A and B channels.
\item
  Reading squared amplitude magnitudes in the A/B basis as channel
  weights.
\item
  A trigonometric parametrization of the amplitudes \(r,t\) as numerical
  implementation coordinates.
\end{enumerate}

Thus, the structure in which the squared magnitude of a complex
amplitude becomes a nonnegative channel weight was already part of the
present model. Theorem 5.1 nevertheless shows that the trigonometric
parametrization need not be an independent derivational assumption: it
can be reconstructed from the projector decomposition of an
exchange-symmetric unitary operator.

\subsubsection{Structures not present in the initial
definitions}\label{structures-not-present-in-the-initial-definitions}

\begin{enumerate}
\def\labelenumi{\arabic{enumi}.}
\tightlist
\item
  Which values of \(R\) produce the special peaks.
\item
  Which phase angles are selected by exact recurrence.
\item
  That the roots are classified by integer pairs \((n,m)\).
\item
  That the selected exchange weights form the discrete sequence
  \(\cos^2(\pi m/n)\).
\item
  That the localization-transfer model and the gray-state metastability
  index connect to finite-order roots of a common linear kernel.
\end{enumerate}

The nontrivial consequence of this work is therefore not the generation
of a square law without assumptions, but rather

\[
\boxed{
\text{finite-order closure}
\Rightarrow
\text{discrete phase sequence}
\Rightarrow
\text{discrete Born-type weight sequence}
}.
\]

\subsection{6.3 Closure-Selected Weights Preceding a Probability
Postulate}\label{63-closure-selected-weights-preceding-a-probability-postulate}

Ordinarily, \(\cos^2\phi\) is introduced as a measurement probability.

Here, the same form first appears as

\begin{quote}
the reflection/transmission weight required for a two-channel exchange
to close exactly after finitely many iterations.
\end{quote}

That is, rather than beginning with

\[
\cos^2 \text{ as a probability},
\]

we first obtain

\[
\cos^2 \text{ as a completely recurrent exchange ratio}.
\]

This order suggests that a Born-type weight may be read not merely as an
external probability postulate, but as a recurrence condition for a
closed phase system.

The word ``selection'' here does not mean dynamical attraction. During
the sweep, \(R\) is not a time-evolving variable; each value of \(R\) is
fixed while the orbital closure error is evaluated. The finite-order
roots are therefore not attractors of the form

\[
R_j\longrightarrow R_{n,m},
\]

but parameters selected by the closure criterion

\[
U_R^n=I.
\]

We call this ``closure selection'' or a ``recurrence-compatible
weight.'' Dynamical selection between A and B under strong observation
is a separate mechanism arising from the D backreaction in System B.

\subsection{6.4 The Discrete Sequence and the Continuous Born
Rule}\label{64-the-discrete-sequence-and-the-continuous-born-rule}

What is obtained directly here is the discrete sequence associated with

\[
\phi_{n,m}=\frac{\pi m}{n}.
\]

As \(n\) increases, rational phases become dense in the continuum of
angles, and hence the set

\[
\cos^2\left(\frac{\pi m}{n}\right)
\]

is dense in \([0,1]\).

Density, however, is not equivalent to having physically derived the
continuous Born rule. To assert a continuous limit as a physical law,
one must separately establish

\begin{itemize}
\tightlist
\item
  which orders are physically realizable,
\item
  how well high-order roots survive noise,
\item
  how observation and coherence times select roots, and
\item
  whether experimental frequencies converge to the squared weights.
\end{itemize}

\begin{center}\rule{0.5\linewidth}{0.5pt}\end{center}

\section{7. Relation to Wave-Packet-Contraction-Like
Behavior}\label{7-relation-to-wave-packet-contraction-like-behavior}

\subsection{7.1 Redistribution of Localization Rather Than Nonunitary
Contraction}\label{71-redistribution-of-localization-rather-than-nonunitary-contraction}

In System A, when one wave packet changes from a weakly localized state
to a strongly localized state, a partial observation can make the packet
appear to contract.

The unnormalized exchange kernel \(U_R\) is unitary; when the linear
update alone is isolated, the total system preserves

\section{\$\$
\textbar A\_j\textbar\^{}2+\textbar B\_j\textbar\^{}2}\label{-a_j2b_j2}

\textbackslash text\{constant\}. \$\$

The full System A update implemented numerically is, however,

\section{\$\$ F\_R(A,B)}\label{-f_rab}

\textbackslash left( \textbackslash mathcal N(rA+tB),
\textbackslash mathcal N(tA+rB) \textbackslash right), \$\$

which normalizes each channel separately. Therefore, \(F_R\) itself
cannot be called a unitary operator.

What occurs at the linear exchange stage is consequently not

\[
\text{contraction accompanied by loss of the total state},
\]

but

\[
\text{redistribution of localization and internal structure within a conserved total state}.
\]

The localization transfer observed numerically in System A is this
conservative redistribution as read after channel-wise normalization.
Whether the finite-order roots are also exact periodic roots of the
complete nonlinear map \(F_R\) does not follow automatically from the
theorem for \(U_R\) and is a System-A-specific proposition requiring
separate verification.

\subsection{7.2 Contraction-Like Appearance Through an Observation
Section}\label{72-contraction-like-appearance-through-an-observation-section}

If an observer reads only one channel or a particular localization index
rather than the entire two-channel state, one section of the combined
orbit of the exchange kernel and normalization may appear as

\begin{itemize}
\tightlist
\item
  a transition from a diffuse to a localized state,
\item
  selection from an intermediate state into A or B, or
\item
  a sudden reduction in wave-packet width.
\end{itemize}

The model therefore offers the working hypothesis

\section{\$\$ \textbackslash boxed\{ \textbackslash text\{apparent
contraction\}}\label{-boxed-textapparent-contraction}

\textbackslash text\{unitary exchange kernel\} +
\textbackslash text\{localization redistribution\} +
\textbackslash text\{channel-wise normalization\} +
\textbackslash text\{restricted observation section\} \}. \$\$

This is not a claim to have reproduced the state-update rule of standard
quantum measurement. It is, however, a concrete mathematical model
capable of producing contraction-like observations without postulating
instantaneous nonunitary contraction as a fundamental operation.

\subsection{7.3 Recurrent Localization Exchange at Finite-Order
Roots}\label{73-recurrent-localization-exchange-at-finite-order-roots}

When the exchange weight is

\[
R=R_{n,m},
\]

the linear kernel \(U_R\) restores every input exactly after \(n\)
iterations.

Thus, even if localization becomes concentrated in one channel at an
intermediate stage of the unnormalized linear evolution, that change is
not irreversible absorption but one phase of a closed orbit. For the
full System A orbit including channel-wise normalization, the conditions
under which this closure survives must be distinguished numerically.

This structure allows the following to be treated as phases of a single
periodic orbit:

\begin{itemize}
\tightlist
\item
  a moment that appears localized,
\item
  a moment that becomes delocalized again,
\item
  a moment at which A and B are interchanged, and
\item
  the moment of exact return to the initial state.
\end{itemize}

\begin{center}\rule{0.5\linewidth}{0.5pt}\end{center}

\section{8. Relation to the White-Cat, Black-Cat, and Gray-Cat
System}\label{8-relation-to-the-white-cat-black-cat-and-gray-cat-system}

\subsection{8.1 A Metastable Two-State
System}\label{81-a-metastable-two-state-system}

In System B, the allocation difference

\[
S_j=|A_j|^2-|B_j|^2
\]

represents the dominance of the two states A and B.

Because the antisymmetric eigenphase rotates, it generally has the form

\[
S_j=C\cos(j\omega+\varphi_0).
\]

The gray state is therefore not merely a static half-and-half mixture of
A and B; it can be represented as a dynamically metastable state in
which the A/B allocation difference is exchanged periodically with small
amplitude.

\subsection{8.2 Weak Readout and Strong-Observation
Selection}\label{82-weak-readout-and-strong-observation-selection}

A one-step update in System B can be written by separating the
unperturbed exchange kernel from the observation maps:

\[
\begin{aligned}
\Psi_{j+1}^{(C)}
&=\mathcal C\!\left(U_R\Psi_j^{(C)}\right),\\
\Psi_{j+1}^{(D)}
&=\mathcal D\!\left(U_R\Psi_j^{(D)}\right).
\end{aligned}
\]

Under weak C readout, the observational action is small and reads the
allocation difference without substantially changing the iterated orbit.
Strong D observation has an explicit nonlinear backreaction that uses
the read value \(S_{D,j}\) to update

\section{\$\$ S\_\{j+1\}}\label{-s_j1-1}

S\_j+g\_D S\_\{D,j\}(1-S\_j\^{}2), \$\$

thereby driving the system toward a state dominated by either A or B.
The direction and irreversibility of A/B selection therefore do not
arise from the finite-order kernel \(U_R\) alone; they depend on the
structure of \(\mathcal D\) and on the coupling strength \(g_D\).

This construction is not a direct implementation of the standard
theories of weak and strong measurement. It nevertheless distinguishes,
within the same model,

\begin{itemize}
\tightlist
\item
  a readout that tends to preserve the state, and
\item
  a selective operation that changes the state.
\end{itemize}

Here, \(U_R\) determines the locations of the finite-order roots,
\(\mathcal C\) reads the gray orbit while disturbing it as little as
possible, and \(\mathcal D\) selects A or B. This division of roles
makes it possible to connect recurrence and observation selection
without confusing them as effects of the same operator.

\subsection{8.3 The Context in Which the Born-Type Weights
Appeared}\label{83-the-context-in-which-the-born-type-weights-appeared}

It is important that the \(\cos^2\) sequence was not discovered from an
abstract two-state vector alone.

It emerged while identifying the cause of peaks in numerical experiments
that evaluated the finite order of the unperturbed exchange kernel in a
system containing

\begin{enumerate}
\def\labelenumi{\arabic{enumi}.}
\tightlist
\item
  two states A and B,
\item
  an intermediate gray metastable state,
\item
  preservation under weak readout,
\item
  A/B selection by a separately implemented strong-observation
  backreaction, and
\item
  closed orbits with small long-time drift.
\end{enumerate}

The principal significance of this work is therefore that, in a family
of mathematical models resembling the physical context in which a
Born-type square law is ordinarily used, the same form of weight emerges
from the finite-order condition of the unperturbed kernel, and the C/D
observation maps can be connected to that same kernel.

\begin{center}\rule{0.5\linewidth}{0.5pt}\end{center}

\section{9. Numerical Peaks and Analytic
Roots}\label{9-numerical-peaks-and-analytic-roots}

\subsection{9.1 Peak Metric}\label{91-peak-metric}

For a time series of length \(k\) in System B, we evaluated

\section{\$\$ \textbackslash overline S\_k}\label{-overline-s_k}

\textbackslash frac1k\textbackslash sum\_\{j=0\}\^{}\{k-1\}S\_j, \$\$

\section{\$\$ A\_k}\label{-a_k}

\textbackslash frac\{\textbackslash max\_\{j\textless k\}S\_j-\textbackslash min\_\{j\textless k\}S\_j\}\{2\},
\$\$

\section{\$\$ D\_k}\label{-d_k}

\subsection{\textbackslash left\textbar{} \textbackslash overline
S\_\{k,\textbackslash mathrm\{second\}\}}\label{left-overline-s_kmathrmsecond}

\textbackslash overline S\_\{k,\textbackslash mathrm\{first\}\}
\textbackslash right\textbar. \$\$

The principal conditions used to confirm the finite-order roots directly
were the initial relative phase

\[
\varphi_0=0\ \text{or}\ \pi,
\]

the code-level initial difference \(s_0=0.01\), and hence the allocation
difference

\[
S_0=|a_0|^2-|b_0|^2=0.02,
\]

with target amplitude \(C=0.02\). Under these conditions, the gray error
was defined as

\section{\$\$ \textbackslash varepsilon\_k}\label{-varepsilon_k}

\textbar\textbackslash overline S\_k\textbar{} +
\textbar A\_k-C\textbar{} + D\_k, \$\$

and the depth as

\[
d_k=-\log_{10}\varepsilon_k.
\]

This metric does not evaluate agreement with a Born probability. It
evaluates

\begin{itemize}
\tightlist
\item
  neutrality of the mean,
\item
  preservation of the oscillation amplitude, and
\item
  absence of drift between the first and second halves.
\end{itemize}

\subsection{9.2 Even-Order Roots}\label{92-even-order-roots}

Under the principal condition \(\varphi_0=0\) or \(\pi\), the allocation
difference at a root of Theorem 5.1 is

\section{\$\$ S\_j}\label{-s_j}

C\textbackslash cos\textbackslash left(\textbackslash frac\{2\textbackslash pi
m\}\{n\}j\textbackslash right), \textbackslash qquad C=0.02. \$\$

If the primitive order \(n\) is even and \(\gcd(m,n)=1\), then \(m\) is
odd, so the lattice points contain exactly

\[
S_0=C,
\qquad
S_{n/2}=C\cos(\pi m)=-C.
\]

The sum of the cyclotomic phases over one period is also zero.

Therefore, when two complete periods \(j=0,\ldots,2n-1\), comprising
\(2n\) samples, are taken, the current gray metric satisfies exactly

\[
\overline S=0,
\qquad
A=C,
\qquad
D=0.
\]

This does not assert that the gray error vanishes for every initial
state and every evaluation window whenever the order is even. Operator
recurrence \(U_R^n=I\) is independent of the initial state, but the
amplitude condition and zero error above are exact results under the
experimental conditions \(\varphi_0=0\) or \(\pi\), \(S_0=C=0.02\), and
a two-complete-period sampling window.

Consequently,

\[
\varepsilon=0,
\]

and in ideal arithmetic

\[
d\rightarrow+\infty.
\]

What diverges is not a physical energy or amplitude but the negative
logarithm of an error metric for a perfectly periodic orbit.
Finite-precision calculations display a finite depth because of rounding
error.

\subsection{9.3 Odd-Order Roots and Observable
Dependence}\label{93-odd-order-roots-and-observable-dependence}

For an odd primitive order \(n\),

\[
U_R^n=I
\]

still holds.

If, however, the finite sampled sequence does not contain both positive
and negative extrema as lattice points, a small deficit remains in the
present amplitude estimate.

Therefore,

\[
\boxed{
\text{exact operator recurrence}
\neq
\text{an infinitely deep peak in a particular observable}
}.
\]

This shows the need to distinguish the candidate sequence of Born-type
weights from how visibly that sequence appears in a particular
experiment.

\subsection{9.4 The Two Observed Principal
Peaks}\label{94-the-two-observed-principal-peaks}

The two principal peaks tracked by the full-range and local
high-precision sweeps coincide with the following roots of Theorem 5.1
{[}4{]}.

{\def\LTcaptype{none} % do not increment counter
\begin{longtable}[]{@{}lrrrr@{}}
\toprule\noalign{}
Root & Primitive order \(n\) & \(m\) & \(R_{n,m}=\lvert r\rvert^2\) &
\(T_{n,m}=\lvert t\rvert^2\) \\
\midrule\noalign{}
\endhead
\bottomrule\noalign{}
\endlastfoot
\(R_{124,23}\) & 124 & 23 & 0.697177927556659 & 0.302822072443341 \\
\(R_{122,23}\) & 122 & 23 & 0.688363946817593 & 0.311636053182407 \\
\end{longtable}
}

Both have even primitive order and odd \(m=23\). Hence, under the
principal conditions of §9.2, the positive and negative extrema occur on
the sampling lattice and the gray error vanishes in ideal arithmetic.
The sharp peaks observed in the numerical sweeps are finite-precision
readouts of this even-order cyclotomic recurrence.

What has been identified here is not the fine-structure constant itself,
but finite-order roots of the exchange operator. The theorem in this
paper does not explain why these roots lie near physical constants; that
remains the separate problem retained in the preceding report {[}4{]}.

\begin{center}\rule{0.5\linewidth}{0.5pt}\end{center}

\section{10. What This Study
Establishes}\label{10-what-this-study-establishes}

\subsection{10.1 Derived Mathematical
Consequences}\label{101-derived-mathematical-consequences}

The following statements are derived rigorously in this work.

\begin{enumerate}
\def\labelenumi{\arabic{enumi}.}
\tightlist
\item
  Up to an overall phase, an exchange-symmetric two-channel unitary
  operator is represented uniquely as
\end{enumerate}

\[
U=P_s+\zeta P_a,
\qquad |\zeta|=1.
\]

\begin{enumerate}
\def\labelenumi{\arabic{enumi}.}
\setcounter{enumi}{1}
\tightlist
\item
  The exchange amplitudes in the channel basis are
\end{enumerate}

\[
r=\frac{1+\zeta}{2},
\qquad
t=\frac{1-\zeta}{2},
\]

so an independent trigonometric assumption is unnecessary.

\begin{enumerate}
\def\labelenumi{\arabic{enumi}.}
\setcounter{enumi}{2}
\item
  After the overall phase has been fixed, the exact recurrence condition
  \(U^n=I\) is equivalent to the antisymmetric eigenphase being an
  \(n\)th root of unity.
\item
  At a nontrivial finite-order root \(\zeta=e^{-2\pi im/n}\), the A/B
  transition weight from an A-basis input is
\end{enumerate}

\[
W_{A\to A}(n,m)=\cos^2\left(\frac{\pi m}{n}\right).
\]

\begin{enumerate}
\def\labelenumi{\arabic{enumi}.}
\setcounter{enumi}{4}
\tightlist
\item
  The complementary weight is
\end{enumerate}

\[
W_{A\to B}(n,m)=\sin^2\left(\frac{\pi m}{n}\right).
\]

\begin{enumerate}
\def\labelenumi{\arabic{enumi}.}
\setcounter{enumi}{5}
\item
  As squared magnitudes of amplitude projections onto the A/B basis,
  these weights use the same mathematical operation and have the same
  \(\cos^2/\sin^2\) form as the standard two-state Born rule.
\item
  This discrete sequence was not obtained by inserting the Born rule as
  the objective function; it follows from exchange symmetry, unitarity,
  and finite-order recurrence.
\item
  The theorem concerns the common linear kernel \(U\) and does not claim
  identity with the complete maps containing channel-wise normalization
  in System A or the C/D backreactions in System B.
\end{enumerate}

\subsection{10.2 Physical Implications Suggested by the
Result}\label{102-physical-implications-suggested-by-the-result}

The study suggests the following possibility:

\begin{quote}
Born-type squared weights can appear as transition weights obtained by
projecting the finitely recurrent eigenphases of an exchange-symmetric
closed two-channel wave system onto its channel basis.
\end{quote}

It further suggests:

\begin{quote}
Localization concentration that appears as wave-packet contraction can
be represented as a composition of a unitary exchange kernel,
redistribution of localization, channel-wise normalization, and a
restricted observation section. Two-state selection under strong
observation can be represented as a separate dynamics obtained by
connecting a nonlinear D backreaction to the same exchange kernel.
\end{quote}

\subsection{10.3 What This Study Does Not Yet
Establish}\label{103-what-this-study-does-not-yet-establish}

This paper does not prove

\begin{enumerate}
\def\labelenumi{\arabic{enumi}.}
\tightlist
\item
  the complete Born rule of standard quantum mechanics,
\item
  a mechanism generating a probabilistic outcome in a single trial,
\item
  convergence of repeated-trial frequencies to \(\cos^2\phi\),
\item
  uniqueness of a general projection measure in arbitrary-dimensional
  Hilbert space,
\item
  irreversibility in an actual quantum measurement,
\item
  complete dynamics of the observer, environment, or decoherence,
\item
  identity between localization transfer in System A and standard
  wave-function collapse,
\item
  exact closure of the complete System A map including channel-wise
  normalization at the same finite order as the linear kernel, or
\item
  unique determination of the same channel weights for a general
  \(U(2)\) operator without exchange symmetry.
\end{enumerate}

Accordingly, this paper does not state that it has ``derived the Born
rule.'' Its claim is instead:

\begin{quote}
Discrete Born-type weights emerge from finite-order recurrence.
\end{quote}

\begin{center}\rule{0.5\linewidth}{0.5pt}\end{center}

\section{11. Relation to the Standard Born
Rule}\label{11-relation-to-the-standard-born-rule}

\subsection{11.1 Born\textquotesingle s Statistical
Interpretation}\label{111-borns-statistical-interpretation}

In 1926, Born introduced the statistical interpretation that reads event
probabilities from the amplitude of the wave function in quantum
scattering {[}5{]}. In modern notation, for a normalized state
\(|\psi\rangle\) and projector \(P_A\),

\[
P(A)=\langle\psi|P_A|\psi\rangle.
\]

For a pure state and a one-dimensional projection,

\[
P(A)=|\langle A|\psi\rangle|^2.
\]

\subsection{11.2 Difference from Gleason\textquotesingle s
Theorem}\label{112-difference-from-gleasons-theorem}

Gleason\textquotesingle s theorem shows that, in Hilbert spaces of
dimension three or greater, an additive probability measure on
orthogonal projections has the density-operator form

\[
\mu(P)=\operatorname{Tr}(\rho P)
\]

{[}6{]}. It is a strong result concerning the mathematical uniqueness of
Born-type measures.

The standard Gleason theorem does not directly include two-dimensional
Hilbert spaces. Busch, however, obtained a generalized Gleason-type
result with a density-operator representation by imposing additivity on
effects, a broader class than projective measurements {[}7{]}. Thus,
established mathematical frameworks also characterize Born-type measures
for two-state systems.

The present study does not address that level of generality. It concerns
an iterated two-channel exchange system and asks

\begin{quote}
Why are particular squared projection values selected as exchange
weights capable of finite recurrence?
\end{quote}

The roles are therefore different:

\begin{itemize}
\tightlist
\item
  Gleason-type results constrain the form of a consistent probability
  measure.
\item
  This work identifies the discrete squared weights selected by the
  closure condition of an iterated exchange dynamics.
\end{itemize}

\subsection{11.3 Difference from Other Derivations of the Born
Rule}\label{113-difference-from-other-derivations-of-the-born-rule}

Zurek\textquotesingle s envariance argument aims to derive
Schmidt-component probabilities \(p_k\propto|\psi_k|^2\) from
entanglement symmetry in a composite system {[}8{]}. It directly
addresses the question of why squared amplitudes should be read as
probabilities.

The question here is different. Without taking probability as the
starting concept, we classify the eigenphases satisfying

\[
U^n=I
\]

for a single exchange-symmetric two-channel operator and show that its
channel-projection weights become \(\cos^2/\sin^2\). This paper is
therefore not an alternative to envariance or Gleason-type measure
theorems, but presents another mathematical route by which Born-type
projection weights arise from finite-order recurrence.

\subsection{11.4 Connection to Exact-Recurrence
Research}\label{114-connection-to-exact-recurrence-research}

Anand et al. classified exact, initial-state-independent recurrence in
finite-dimensional Floquet systems through the cyclotomic structure of
unitary spectra {[}9{]}. The condition

\[
\zeta^n=1
\]

in this paper is the minimal two-channel realization of such exact
recurrence. Whereas that preceding work gives an arithmetic
classification of recurrence times for general Floquet spectra, exchange
symmetry here fixes the eigenvectors to symmetric and antisymmetric
modes. Returning their cyclotomic eigenphases to the A/B channel basis
yields

\section{\$\$ \textbackslash left\textbar{}
\textbackslash frac\{1\textbackslash pm\textbackslash zeta\}\{2\}
\textbackslash right\textbar\^{}2}\label{-left-frac1pmzeta2-right2}

\textbackslash cos\^{}2\textbackslash phi,\textbackslash{}
\textbackslash sin\^{}2\textbackslash phi. \$\$

This mapping from cyclotomic recurrence to channel squared weights is
the particular focus of the present paper.

\subsection{11.5 Significance of Restricting the Analysis to a Two-State
System}\label{115-significance-of-restricting-the-analysis-to-a-two-state-system}

Because this is a two-state system, the Born rule is not derived by
directly applying Gleason\textquotesingle s theorem.

Two-state systems nevertheless form the minimal models of

\begin{itemize}
\tightlist
\item
  interference,
\item
  two-level spin,
\item
  two-path measurement,
\item
  two-mode optics, and
\item
  A/B observation selection.
\end{itemize}

Showing that recurrent exchange weights in this minimal model appear as

\[
\cos^2\phi,
\qquad
\sin^2\phi
\]

is a result with a clearly delimited domain of validity before any
extension to a general theory.

\begin{center}\rule{0.5\linewidth}{0.5pt}\end{center}

\section{12. Falsifiable Tests}\label{12-falsifiable-tests}

\subsection{\texorpdfstring{12.1 The Exchange-Symmetric Subgroup and
General
\(U(2)\)}{12.1 The Exchange-Symmetric Subgroup and General U(2)}}\label{121-the-exchange-symmetric-subgroup-and-general-u2}

Theorem 5.1 rules out the possibility that the \(\cos^2\) form arises
only from the scattering-coefficient representation of §3. What is
required is not a particular trigonometric coordinate but the A/B
exchange symmetry

\[
[U,X]=0.
\]

A general two-level unitary operator, including its overall phase, can
be written as

\section{\$\$ U}\label{-u}

e\^{}\{i\textbackslash alpha\} \textbackslash left(
\textbackslash cos\textbackslash beta,I
-i\textbackslash sin\textbackslash beta,\textbackslash boldsymbol
n\textbackslash cdot\textbackslash boldsymbol\textbackslash sigma
\textbackslash right), \textbackslash qquad
\textbar\textbackslash boldsymbol n\textbar=1. \$\$

The finite-order condition quantizes the eigenphase \(\beta\) to a
rational angle, but the transition weight from the A basis to the B
basis is

\section{\$\$ \textbar\textbackslash langle
B\textbar U\textbar A\textbackslash rangle\textbar\^{}2}\label{-langle-buarangle2}

\textbackslash sin\^{}2\textbackslash beta,(n\_x\^{}2+n\_y\^{}2), \$\$

which retains the eigenvector axis \(\boldsymbol n\). Finite order alone
therefore cannot uniquely fix the general-\(U(2)\) A/B weight to
\(\sin^2\beta\).

The theorem here applies to the centralizer of the exchange operator
\(X=\sigma_x\): the case in which \(\boldsymbol n\) is fixed to the
exchange axis and the symmetric and antisymmetric eigenvectors have
equal weights relative to the A/B basis. This is not a defect of the
theorem but the identification of the precise symmetry condition that
produces the discrete Born-type weights.

A falsifiable next test is to introduce a controlled perturbation with
\([U,X]\ne0\) and measure whether the peak weights change according to
the axis factor \(n_x^2+n_y^2\). If they do not, an additional structure
not explained by the projection mechanism in this paper must be present.

\subsection{12.2 Testing the Frequency
Law}\label{122-testing-the-frequency-law}

The present result gives closed-orbit weights but does not directly
generate measurement frequencies.

The next stage must prepare the same A-basis input \(|A\rangle\)
repeatedly, pass it through the finite-order kernel and an explicit
measurement/selection map, and measure the A/B output frequencies

\[
f_{A\to A},
\qquad
f_{A\to B},
\]

testing whether

\[
f_{A\to A}\rightarrow W_{A\to A}(n,m)=|r_{n,m}|^2,
\qquad
f_{A\to B}\rightarrow W_{A\to B}(n,m)=|t_{n,m}|^2.
\]

If this holds, it will provide a more direct connection between
finite-order exchange weights and the Born rule as an experimental
frequency law. For a general superposed input, a frequency law
containing the interference term in \(r\alpha+t\beta\) must be tested
separately.

\subsection{12.3 Observational Perturbation and
Irreversibility}\label{123-observational-perturbation-and-irreversibility}

A fully unitary system is recurrent and is not irreversible in
principle.

To connect the model to actual observational selection, one must
introduce

\begin{itemize}
\tightlist
\item
  apparatus degrees of freedom,
\item
  environmental degrees of freedom,
\item
  phase noise,
\item
  channel loss, and
\item
  post-readout feedback,
\end{itemize}

and determine the conditions under which a recurrent orbit becomes
effectively fixed in either A or B.

\subsection{12.4 Finite-Order Selection Under
Noise}\label{124-finite-order-selection-under-noise}

Let the displacement from a root be

\[
R=R_{n,m}+\delta R.
\]

The resonance width of each finite-order root can be defined by
measuring the operator residual after \(n\) iterations,

\[
\mathcal E_n(R)
:=
\|U_R^n-I\|.
\]

One may also introduce per-iteration phase noise

\[
\theta_j=\theta+\xi_j
\]

and evaluate the recurrence fidelity of

\[
U_{j+n-1}\cdots U_j.
\]

This will distinguish which low-order roots among the mathematically
dense set are physically robust.

\subsection{12.5 Normalization Audit of System
A}\label{125-normalization-audit-of-system-a}

The System A implementation contains channel-wise normalization, so its
complete update map is not the linear operator \(U_R\otimes I\) but a
state-dependent nonlinear map.

It is therefore necessary to compare three conditions:

\begin{enumerate}
\def\labelenumi{\arabic{enumi}.}
\tightlist
\item
  no normalization,
\item
  joint normalization of the entire two-channel state, and
\item
  separate normalization of each channel.
\end{enumerate}

This comparison must separate the location of the finite-order roots,
localization transfer, and the visibility of Born-type weights.

Such an audit is required to connect rigorously the finite-order
structure derived for the common linear kernel to the nonlinear
wave-packet localization behavior of System A. It has already been
confirmed that the current implementation uses channel-wise
normalization. What remains unconfirmed is under which conditions that
normalization preserves, shifts, or eliminates the finite-order roots of
the linear kernel.

\begin{center}\rule{0.5\linewidth}{0.5pt}\end{center}

\section{13. Discussion}\label{13-discussion}

\subsection{13.1 Principal Insight}\label{131-principal-insight}

The principal insight of this study is that a seemingly continuous wave
system with a continuous exchange coefficient \(R\) contains exact
discrete closed orbits classified by integer pairs \((n,m)\), and that
their weights are fixed directly by exchange symmetry.

Moreover, the exchange weights of these discrete closed orbits are

\[
R_{n,m}=\cos^2\left(\frac{\pi m}{n}\right).
\]

Exchange symmetry and unitarity give

\[
U=P_s+\zeta P_a,
\]

while finite-order closure gives \(\zeta=e^{-2\pi im/n}\). Hence the
structure

\[
\text{discrete phase closure}
\longrightarrow
\text{amplitude projection}
\longrightarrow
\text{squared weight}
\]

follows. The \(\cos^2\) form is not a remnant of an assumed scattering
coordinate; it is the squared interference obtained when the sum and
difference of two projected eigenphases are returned to the A/B basis.

\subsection{13.2 A Finite-Order Origin of Born-Type
Weights}\label{132-a-finite-order-origin-of-born-type-weights}

In the Born rule, squared weights are given as observation
probabilities.

Here, the same squared weights appear as the exchange ratios required to
return to an identical state after finitely many iterations.

The following proposition is therefore a theorem within this model:

\begin{quote}
In an exchange-symmetric closed two-channel system, transition weights
obtained by projecting recurrent cyclotomic eigenphases onto the A/B
channels form a discrete Born-type squared sequence.
\end{quote}

In this derivation, probability is not posited initially as an
unstructured random number. The order is instead

\begin{enumerate}
\def\labelenumi{\arabic{enumi}.}
\tightlist
\item
  a closed phase relation,
\item
  an integer ratio permitting finite recurrence,
\item
  projection onto the observation basis, and
\item
  squaring of the projected component.
\end{enumerate}

What remains unresolved is the mechanism by which these transition
weights govern physical repeated-measurement frequencies, not the
derivation of the discrete squared weights themselves.

\subsection{13.3 A Common Kernel for Contraction, Metastability, and
Selection}\label{133-a-common-kernel-for-contraction-metastability-and-selection}

Systems A and B were constructed to investigate different phenomena.

After analysis, however,

\begin{itemize}
\tightlist
\item
  concentration of localization that appears as wave-packet contraction,
\item
  periodic exchange between the two states A and B,
\item
  a gray metastable state,
\item
  preservation under weak readout,
\item
  selection under strong observation, and
\item
  discrete Born-type weights
\end{itemize}

were all connected to an unperturbed kernel having the same
decomposition into symmetric and antisymmetric eigenmodes. The complete
update maps are nevertheless not identical. The localization
representation in System A adds channel-wise normalization
\(\mathcal N_A\); weak readout in System B adds \(\mathcal C\); and
strong-observation selection adds nonlinear backreaction \(\mathcal D\).

This means that phenomena discussed as separate problems in quantum
measurement can be compared, in the minimal two-channel model, as
combinations of a ``common exchange-phase kernel'' and distinct
downstream maps. The important point is that the root location,
localization representation, weak readout, and strong selection can be
connected in the same computational model without conflating them as a
single effect.

\subsection{13.4 Exact Central
Conclusion}\label{134-exact-central-conclusion}

The exact central conclusion of this paper is:

\begin{quote}
The common linear kernel of model families independently constructed to
study wave-packet localization transfer, metastable two-state dynamics,
weak readout, and strong-observation selection is an exchange-symmetric
two-channel unitary operator. Finite-order closure of this kernel
necessarily yields the discrete channel-projection weights
\(\cos^2(\pi m/n)\) and \(\sin^2(\pi m/n)\) without specifying the Born
rule as the target of the search.
\end{quote}

This is not a derivation of the complete Born rule. It does show, by a
reproducible model and analytic formulas, that discrete Born-type
squared weights arise from exchange symmetry, unitarity, finite-order
closure, and channel projection. The frequency law, nonlinear closure of
System A, and correspondence between D observation and standard
measurement dynamics are separate next problems, not parts of this
theorem.

\begin{center}\rule{0.5\linewidth}{0.5pt}\end{center}

\section{14. Conclusion}\label{14-conclusion}

This study analyzed two-channel unitary operators commuting with the
channel-exchange operator \(X\). By exchange symmetry, after fixing the
overall phase the operator can be written as

\[
U=P_s+\zeta P_a,
\qquad |\zeta|=1.
\]

Returning to the channel basis gives the amplitudes

\[
r=\frac{1+\zeta}{2},
\qquad
t=\frac{1-\zeta}{2}.
\]

The finite-order condition \(U^n=I\) quantizes the antisymmetric
eigenphase to

\[
\zeta=e^{-2\pi im/n},
\]

and yields the finitely recurrent channel-projection weight

\[
\boxed{
|r_{n,m}|^2=\cos^2\left(\frac{\pi m}{n}\right)
}
\]

and its complement

\[
\boxed{
|t_{n,m}|^2=\sin^2\left(\frac{\pi m}{n}\right)
}.
\]

As squared magnitudes of the amplitudes obtained by projecting an
A-basis input onto the A/B basis, this \(\cos^2/\sin^2\) form has the
same mathematical structure as the standard two-state Born rule in
quantum mechanics. The trigonometric form is not an independent
assumption: it follows from interference between the exchange-symmetric
projectors \(P_s,P_a\) and a cyclotomic eigenphase.

The important point is that the model was not designed to reproduce this
form. It began with two numerical models intended to study wave-packet
localization transfer, contraction-like localization, the
white-cat/black-cat/gray-cat metastable interface, weak readout, and
strong-observation selection. Subsequent eigenvalue analysis of the
observed sharp peaks identified the finite-order roots and discrete
Born-type weights, which this paper then elevated to an
exchange-symmetric finite-order theorem independent of the trigonometric
parametrization.

The roles of the implementation have also been separated. The common
linear kernel \(U\) determines the root locations;
\(\mathcal N_A\circ U\), including channel-wise normalization, produces
the localization representation of System A; and \(\mathcal C\circ U\)
and \(\mathcal D\circ U\) produce weak readout and strong selection,
respectively, in System B.

The conclusion of this paper is therefore not a complete derivation of
the Born rule.

What has been established is the precise mathematical connection

\[
\boxed{
\begin{gathered}
\text{exchange-symmetric unitary kernel}
+\text{cyclotomic finite-order closure}\\
+\text{A/B channel projection}
\Rightarrow\text{discrete Born-type squared weights}
\end{gathered}
}.
\]

The next stage will examine general \(U(2)\) control experiments that
break exchange symmetry, reproduction of measurement frequencies,
inclusion of apparatus and environmental degrees of freedom, selection
of roots under noise, and an audit of System A under different
normalization conditions.

\begin{center}\rule{0.5\linewidth}{0.5pt}\end{center}

\section{References}\label{references}

\subsection{Self-Citations}\label{self-citations}

\begin{enumerate}
\def\labelenumi{\arabic{enumi}.}
\tightlist
\item
  Noriaki Kihara, ``Experimental Specification v1 for Low-Localization
  and Harmonic Transfer Readout in Fermion-Like Collisions Using an
  Exchange-Interference Scattering Matrix,'' 2026. {[}Japanese{]}.
\item
  Noriaki Kihara, ``Preliminary Experimental Summary v1 of the
  Acceleration Basis and Localization Exchange in Fermion-Like
  Collisions Using an Exchange-Interference Scattering Matrix,'' Zenodo
  Concept DOI: 10.5281/zenodo.21333766, 2026. {[}Japanese{]}.
\item
  Noriaki Kihara, ``Preliminary Experimental Summary v1 of C Weak
  Readout and D Strong-Observation Selection at a White-Cat, Black-Cat,
  and Gray-Cat Metastable Interface,'' Zenodo Concept DOI:
  10.5281/zenodo.21353208, 2026. {[}Japanese{]}.
\item
  Noriaki Kihara, ``Discovery of Finite-Order Resonances in Iterated
  Exchange Scattering: Identification of the Origin of Peaks Near the
  Fine-Structure Values 137 and 128 and a Reproducible Wave-Packet
  Mathematical Model,'' Version DOI: 10.5281/zenodo.21421367, Concept
  DOI: 10.5281/zenodo.21421366, 2026.
\end{enumerate}

\subsection{External References}\label{external-references}

\begin{enumerate}
\def\labelenumi{\arabic{enumi}.}
\setcounter{enumi}{4}
\tightlist
\item
  Max Born, ``Zur Quantenmechanik der Stoßvorgänge,'' \emph{Zeitschrift
  für Physik} \textbf{37}, 863--867 (1926). DOI: 10.1007/BF01397477.
\item
  Andrew M. Gleason, ``Measures on the Closed Subspaces of a Hilbert
  Space,'' \emph{Journal of Mathematics and Mechanics} \textbf{6},
  885--893 (1957). DOI: 10.1512/iumj.1957.6.56050.
\item
  Paul Busch, ``Quantum States and Generalized Observables: A Simple
  Proof of Gleason\textquotesingle s Theorem,'' \emph{Physical Review
  Letters} \textbf{91}, 120403 (2003). DOI:
  10.1103/PhysRevLett.91.120403.
\item
  Wojciech H. Zurek, ``Probabilities from Entanglement:
  Born\textquotesingle s Rule \(p_k=|\psi_k|^2\) from Envariance,''
  \emph{Physical Review A} \textbf{71}, 052105 (2005). DOI:
  10.1103/PhysRevA.71.052105.
\item
  Amit Anand, Dinesh Valluri, Jack Davis, and Shohini Ghose, ``Quantum
  Recurrences and the Arithmetic of Floquet Dynamics,'' \emph{Quantum}
  \textbf{10}, 2074 (2026). DOI: 10.22331/q-2026-04-20-2074.
\end{enumerate}

\begin{center}\rule{0.5\linewidth}{0.5pt}\end{center}

\section{Appendix A. Classification of
Claims}\label{appendix-a-classification-of-claims}

{\def\LTcaptype{none} % do not increment counter
\small
\begin{longtable}[]{@{}p{0.54\linewidth}p{0.39\linewidth}@{}}
\toprule\noalign{}
Claim & Classification \\
\midrule\noalign{}
\endhead
\bottomrule\noalign{}
\endlastfoot
Phase-fixed representation \(U=P_s+\zeta P_a\) of an exchange-symmetric
unitary operator & Derived consequence of Theorem 5.1 \\
Finite-order condition \(U^n=I\Leftrightarrow\zeta^n=1\) & Derived
consequence of Theorem 5.1 \\
\(r=(1+\zeta)/2\), \(t=(1-\zeta)/2\) & Derived from the projector
decomposition \\
\(\lvert r_{n,m}\rvert^2=\cos^2(\pi m/n)\) & Derived consequence \\
\(\lvert t_{n,m}\rvert^2=\sin^2(\pi m/n)\) & Derived consequence \\
Trigonometric form is recovered as an implementation coordinate rather
than an independent assumption & Derived consequence \\
Isomorphism between finite-order weights for a basis input and two-state
Born projection & Derived comparison \\
The \(\cos^2\) sequence was not entered as the search target & Fact
about experimental and development history \\
Systems A and B share the same unnormalized, unperturbed linear kernel &
Implementation fact confirmed in source code \\
System A uses channel-wise normalization and its complete map is
nonlinear & Implementation fact confirmed in source code \\
The complete nonlinear map of System A recurs exactly with the same
order as the linear kernel & Untested \\
D selection in System B has an explicit nonlinear backreaction &
Implementation fact confirmed in source code \\
The gray error vanishes at even-order roots & Derived consequence
restricted to \(\varphi_0=0\) or \(\pi\), \(S_0=C=0.02\), and \(2n\)
samples \\
Infinitely deep peak & Mathematical divergence of the error metric
\(-\log_{10}\varepsilon\); not an energy divergence \\
Finite order alone does not determine the same A/B weights for general
\(U(2)\) & Derived consequence of §12.1 \\
Contraction-like wave-packet behavior can be represented by the exchange
kernel, normalization, and partial observation & Mathematical finding
within the present System A model \\
Standard quantum-mechanical wave-function collapse has been explained &
Not derived \\
The complete Born rule has been derived & Not derived \\
Repeated-measurement frequencies converge to \(\cos^2\) & Untested \\
Discrete Born-type channel weights originate from finite-order phase
closure & Derived for the exchange-symmetric two-channel kernel \\
The physical Born frequency law originates from finite-order phase
closure & Physical hypothesis for future testing \\
\end{longtable}
}

\begin{center}\rule{0.5\linewidth}{0.5pt}\end{center}

\section{Appendix B. Minimal Reproduction
Formulas}\label{appendix-b-minimal-reproduction-formulas}

The central result of this paper can be reproduced without assuming a
trigonometric representation, using only the following formulas.

\subsection{B.1 Exchange-Symmetric
Projectors}\label{b1-exchange-symmetric-projectors}

\[
X=
\begin{pmatrix}
0&1\\
1&0
\end{pmatrix},
\]

\[
P_s=\frac{I+X}{2},
\qquad
P_a=\frac{I-X}{2}.
\]

\subsection{B.2 Exchange-Symmetric Unitary
Operator}\label{b2-exchange-symmetric-unitary-operator}

If

\[
[U,X]=0,
\qquad
U^\dagger U=I,
\]

then, after fixing the overall phase,

\[
\boxed{
U=P_s+\zeta P_a,
\qquad |\zeta|=1
}
\]

can be written.

\subsection{B.3 Finite-Order Condition}\label{b3-finite-order-condition}

For a nontrivial primitive order \(n\ge3\),

\[
U^n=I
\quad\Longleftrightarrow\quad
\zeta^n=1.
\]

Thus, with \(\gcd(m,n)=1\),

\[
\zeta=e^{-2\pi im/n}.
\]

\subsection{B.4 Channel Amplitudes and Squared
Weights}\label{b4-channel-amplitudes-and-squared-weights}

\section{\$\$ U}\label{-u-1}

\section{\textbackslash frac12 \textbackslash begin\{pmatrix\}
1+\textbackslash zeta\&1-\textbackslash zeta\textbackslash{}
1-\textbackslash zeta\&1+\textbackslash zeta
\textbackslash end\{pmatrix\}}\label{frac12-beginpmatrix-1zeta1-zeta-1-zeta1zeta-endpmatrix-1}

\textbackslash begin\{pmatrix\} r\&t\textbackslash{} t\&r
\textbackslash end\{pmatrix\}, \$\$

and hence

\[
r=\frac{1+\zeta}{2},
\qquad
t=\frac{1-\zeta}{2}.
\]

Their squared magnitudes are

\[
\boxed{
|r_{n,m}|^2
=\cos^2\left(\frac{\pi m}{n}\right)
}
\]

and

\[
\boxed{
|t_{n,m}|^2
=\sin^2\left(\frac{\pi m}{n}\right)
}.
\]

These are the discrete Born-type weights of an exchange-symmetric
two-channel process that recurs exactly after finitely many iterations.

\subsection{B.5 Recovery of the Numerical Implementation
Coordinates}\label{b5-recovery-of-the-numerical-implementation-coordinates}

The implementation parametrization

\[
t=e^{i\theta}\cos\theta,
\qquad
r=-ie^{i\theta}\sin\theta
\]

is a coordinate representation of the same operator obtained by setting

\[
\zeta=-e^{2i\theta}.
\]

Indeed, the sum and difference in B.4 give

\section{\$\$
\textbackslash frac\{1+\textbackslash zeta\}\{2\}}\label{-frac1zeta2}

\section{\textbackslash frac\{1-e\^{}\{2i\textbackslash theta\}\}\{2\}}\label{frac1-e2itheta2}

-ie\^{}\{i\textbackslash theta\}\textbackslash sin\textbackslash theta=r,
\$\$

\section{\$\$
\textbackslash frac\{1-\textbackslash zeta\}\{2\}}\label{-frac1-zeta2}

\section{\textbackslash frac\{1+e\^{}\{2i\textbackslash theta\}\}\{2\}}\label{frac1e2itheta2}

e\^{}\{i\textbackslash theta\}\textbackslash cos\textbackslash theta=t.
\$\$

At a finite-order root,

\section{\$\$ \textbackslash theta}\label{-theta-1}

\textbackslash frac\{\textbackslash pi\}\{2\}-\textbackslash frac\{\textbackslash pi
m\}\{n\} \textbackslash pmod\textbackslash pi, \$\$

which reproduces the weights in B.4 directly.

\begin{center}\rule{0.5\linewidth}{0.5pt}\end{center}

\section{Appendix C. Correspondence with the Core
Implementations}\label{appendix-c-correspondence-with-the-core-implementations}

The core updates are shown here so that readers can verify both that the
finite-order roots were not inserted into the code as hidden constants
and that the linear kernel is separated from the downstream maps.

\subsection{C.1 System A: Channel-Wise Normalization After Linear
Exchange}\label{c1-system-a-channel-wise-normalization-after-linear-exchange}

The iterated update in the
\href{../20260715/run_system_A_localization_exchange_R_sweep_preliminary_v1.py}{System
A implementation} is

\begin{Shaded}
\begin{Highlighting}[]
\NormalTok{a\_next }\OperatorTok{=}\NormalTok{ src.normalize(r }\OperatorTok{*}\NormalTok{ a }\OperatorTok{+}\NormalTok{ t }\OperatorTok{*}\NormalTok{ b)}
\NormalTok{b\_next }\OperatorTok{=}\NormalTok{ src.normalize(t }\OperatorTok{*}\NormalTok{ a }\OperatorTok{+}\NormalTok{ r }\OperatorTok{*}\NormalTok{ b)}
\NormalTok{a, b }\OperatorTok{=}\NormalTok{ a\_next, b\_next}
\end{Highlighting}
\end{Shaded}

This explicitly contains the common linear kernel

\[
(a,b)\mapsto(ra+tb,\ ta+rb)
\]

and separate normalization of each output channel. The complete
localization-transfer map is therefore nonlinear.

\subsection{C.2 System B: Unperturbed Finite-Order
Kernel}\label{c2-system-b-unperturbed-finite-order-kernel}

The iterated update in the
\href{../20260715/run_minimal_system_B_gray_direct_check_v5.py}{System B
direct finite-order test} is

\begin{Shaded}
\begin{Highlighting}[]
\NormalTok{a, b }\OperatorTok{=}\NormalTok{ normalize\_pair(r }\OperatorTok{*}\NormalTok{ a }\OperatorTok{+}\NormalTok{ t }\OperatorTok{*}\NormalTok{ b, t }\OperatorTok{*}\NormalTok{ a }\OperatorTok{+}\NormalTok{ r }\OperatorTok{*}\NormalTok{ b)}
\end{Highlighting}
\end{Shaded}

As long as \(U_R\) is unitary and the input pair is normalized, the
normalization factor in \texttt{normalize\_pair} is one. Thus, the
iteration that determines the peak locations in this direct test is
exactly the linear kernel \(U_R\); neither \(R_{124,23}\) nor
\(R_{122,23}\) is injected into the update equation as a target value.

\subsection{C.3 System B: Nonlinear Backreaction of Strong D
Observation}\label{c3-system-b-nonlinear-backreaction-of-strong-d-observation}

The
\href{../20260714/run_gray_cat_d_observation_response_preliminary_v1.py}{strong-D-observation
implementation} adds the following backreaction after exchange:

\begin{Shaded}
\begin{Highlighting}[]
\NormalTok{s\_next }\OperatorTok{=}\NormalTok{ s }\OperatorTok{+}\NormalTok{ d\_gain }\OperatorTok{*}\NormalTok{ s\_d }\OperatorTok{*}\NormalTok{ (}\FloatTok{1.0} \OperatorTok{{-}}\NormalTok{ s }\OperatorTok{*}\NormalTok{ s)}
\NormalTok{a\_next }\OperatorTok{=}\NormalTok{ math.sqrt(}\FloatTok{0.5} \OperatorTok{*}\NormalTok{ (}\FloatTok{1.0} \OperatorTok{+}\NormalTok{ s\_next)) }\OperatorTok{*}\NormalTok{ phase\_a}
\NormalTok{b\_next }\OperatorTok{=}\NormalTok{ math.sqrt(}\FloatTok{0.5} \OperatorTok{*}\NormalTok{ (}\FloatTok{1.0} \OperatorTok{{-}}\NormalTok{ s\_next)) }\OperatorTok{*}\NormalTok{ phase\_b}
\end{Highlighting}
\end{Shaded}

The finite-order roots therefore arise from the cyclotomic spectrum of
\(U_R\), while strong selection into A or B arises from a separate D
backreaction. The code structure itself implements the
paper\textquotesingle s two-layer distinction

\[
\boxed{
\text{generation of recurrence roots}
\ne
\text{state selection under strong observation}
}.
\]

\end{document}
